

\section{Seventeenth Sunday after Trinity: The Greatest in the Kingdom of Heaven}

\begin{quote}

Matthew 18:1–7. At the same time the disciples came to Jesus and said, “Who is the greatest in the Kingdom of Heaven?” And Jesus called a little child to Him and set it in the midst of them, and said, “Truly I say to you: Unless you are converted and become as children, you shall by no means enter the Kingdom of Heaven. Therefore, whoever humbles himself as this little child, he is the greatest in the Kingdom of Heaven. And whoever receives such a child in My name receives Me. But whoever causes one of these little ones, who believe in Me, to stumble, it would be better for him if a millstone were hung about his neck, and he were sunk into the depth of the sea. Woe to the world because of stumbling-blocks! For it is necessary that stumbling-blocks come; yet woe to that man by whom the stumbling-block comes!”

\end{quote}

\bigskip

This event falls in that period of Jesus’ earthly life when for the last time He journeyed up from Galilee to Jerusalem. He was going to suffering, cross, and death; but His disciples did not understand it; they were thinking of the Kingdom and the glory, the Kingdom of God, the Kingdom of Heaven, the Kingdom of the Messiah and His glory. Therefore, while Jesus went to become the lowest and most despised and humiliated among the children of men, His disciples in the folly of their minds asked Him, “Who is the greatest in the Kingdom of Heaven?”

There is so dreadful a contrast between the thoughts that occupied the Savior and the thoughts that stirred in the hearts of the disciples, that we must marvel at the heavenly wisdom and love with which our Savior meets their foolish question.

Yet it is a thorough and serious reproof that the disciples receive. A little child is the instrument that the Lord employs in order to rebuke and shame the high-striving disciples.

The first thing Jesus says about the child is a piercing warning: “Truly I say to you: Unless you are converted and become as children, you shall by no means enter the Kingdom of Heaven.” The disciples were thinking about who was greatest in the Kingdom of Heaven; but the Lord wills that first and foremost they should think about how they may enter the Kingdom of Heaven at all. The disciples surely thought that in that matter all was in order; if they were not in the Kingdom of Heaven, who then should be there? Why, they had left all things and followed Jesus; they had believed and confessed that He was Christ, the Son of the living God; they were His friends, His only followers; they must be as good as assured of being in the Kingdom of Heaven; for them the question was: Who is the greatest?

But Jesus says to them, and to all of us who seem so sure that we are within the Kingdom of Heaven: “Unless you are converted and become as children, you shall by no means enter the Kingdom of Heaven.” There is no other way, then, than this: you must become as a child and continue to be as a child, if you would enter in. But that means that it is so difficult to enter in that there is no time or occasion to ask who shall be greatest. If only we might enter in; if only we might be saved! If only we could attain the very lowest place in the Kingdom of God, O how blessed, O how well helped we would be!

Is it then so difficult? Yes, it is difficult to be converted and become as children. To descend from all one’s own height and all one’s own wisdom, gained through a long and painful human life, to cast it away as useless, indeed as a harmful burden that only hinders the soul on its heavenward way, O, that is hard, dreadfully hard.

To have to confess that all our way in the world and every step we have taken has only been vain and empty, that is pain, deadly pain to our natural mind. The great, proud, grown mind, so worldly-wise and so heartless, looks down with contempt upon the childlike mind, which believes so easily and loves so warmly. And yet you must return to that, my proud, suspicious, self-righteous heart; for without this spirit you shall by no means enter the Kingdom of Heaven.

But I cannot, says the grown person. And he speaks the truth. For this is the work of the Holy Spirit. God’s Spirit must work this, that an old human being becomes a child again and may begin afresh with childlike faith and love. And as painful as it is to have to turn back when one thinks one has come so far forward, so blessed is it also to feel childlike peace and childlike joy in one’s old heart, where the fire of the world has burned and the storm-wind of the world has played with the ashes. It is the greatest pain and the highest joy of human life that meet in a mysterious moment when the Spirit from God conquers the spirit of the world in the human soul and gives peace there through faith in Jesus Christ.

Are you then afraid to pass through the pain? Is it so hard to lose the world in order to gain the Kingdom of God? Yet it is altogether harder one day, in the hour of death, to lose the world and to gain an eternal perdition. Continue until death with your worldly wisdom and your worldly desire, and you lose all without gaining anything but eternal torment. Turn back in time, while the Savior calls you to His heart and His cross, and you gain the Kingdom of God with righteousness, peace, and joy. Which do you choose?

But the Savior wills that the child He has set among us shall teach us still more; therefore He adds: “Whoever humbles himself as this little child, he is the greatest in the Kingdom of Heaven.”

It is possible to become great in the Kingdom of Heaven. The Savior Himself shows the way there: “Go and sit down in the lowest place!” If you truly do this, if you do it because you know yourself to be the very least and the most unworthy, then one day there shall sound to you a heavenly, blessed voice: “Friend, move up higher!”

What are you, other than a poor sinner before God? What is there in the oldest and most tried Christian that is good, except that which he has received by grace, nothing but grace? There is nothing else that saves him than grace, grace alone. Why should he shrink from being the least and from sitting down in the lowest place? And if the greatest and the chief became the servant of all, what harm would have been done? Has not the Lord of glory become the servant of all with His own blood, with His shameful death? And to follow in His footsteps and be misunderstood and despised and mocked and scorned as He was, that you will surely not refuse, Christian soul?

Only if you have the childlike spirit, which neither wonders nor is offended that the world passes you by so proudly, then the heart does not become bitter in suffering, nor evil under scorn. Believe in Him who bore the cross before you, and bore the cross for your sake, and His way to glory shall become your way to eternal blessedness. The one who believes most and loves best, he shall become greatest, even though the world should tread him down with the coldest contempt.

But the one who causes God’s children to stumble, who mocks their faith and thrusts away their love, who lays stones in their way so that they may stumble and fall, his own way is dark and his work evil. And yet it is unavoidable that the world will attempt this work of wickedness also. With its cleverness it will try to shoot the arrows of doubt into the hearts of the children; with its desire it will try to lay snares for their feet; with its power it will try to awaken fear and anguish in their souls. Evil, cruel world; it has no peace itself, and those who have peace it would make peace-less again. It shall receive its judgment; but you, child of God, fear the power of seduction and of stumbling; for it is great. Flee to your Savior’s heart and cling as a little child to His hand, and He shall lead you past the stumbling-block, so that it shall not seize you. Safe in your Savior’s arms, with a childlike spirit you shall come well through; and glad to attain the lowest place in the Kingdom of God, you shall learn that the one who needed most grace and therefore clung most closely to the Savior’s heart found the best place. Blessed is the one who became the least, when the least son or daughter is granted to rest nearest the Father’s breast. Then you shall be content with your lot and praise that God who resists the proud, but gives grace to the humble.

